% Separate post-CP35 candidate; no frozen canonical or build input is changed.
\subsection{A primitive globally generated polarization of type \((1,3)\)}

Work over an algebraically closed field \(k\) of characteristic two.
Relative Frobenius on an abelian surface is denoted
\(F:A\to A^{(2)}\). Absolute Frobenius on a coefficient space
will be written \(F_X\). These are different maps.
The full-image assertion below assumes that the ample line bundle
of type \((1,3)\) is globally generated.

\begin{lemma}[Odd Euler characteristic gives Frobenius stability]
\label{lem:DE-odd-line-Frobenius-stable}
\status{Proved; full independent proof, source and code audits passed.}
If \(L\) is a line bundle on an abelian surface and \(\chi(L)\)
is odd, the rank-four vector bundle \(F_*L\) is slope stable
with respect to every ample polarization on \(A^{(2)}\).
\end{lemma}
\begin{proof}
The canonical Frobenius filtration gives \(F^*F_*L\) graded
pieces \(L,L\otimes\Omega_A^1,L\otimes\bigwedge^2\Omega_A^1\).
This is the case \(p=n=2\) of the higher-dimensional filtration
in \citep[Definition~3.1, Lemmas~3.2--3.5 and Theorem~3.6,
pp.~5--9]{Sun2006FrobeniusFiltration}.
Here it also follows directly from the diagonal ideal on
\(A\times_{A^{(2)}}A\): in local smooth coordinates it is
generated by two elements \(\epsilon_1,\epsilon_2\), with
\(\epsilon_1^2=\epsilon_2^2=0\), and its successive degrees
have ranks \(1,2,1\). Their changes of coordinates are those
of \(\mathcal O_A,\Omega_A^1,\bigwedge^2\Omega_A^1\).
Since \(\Omega_A^1\simeq\mathcal O_A^2\), the bundle
\[
                U=F^*F_*L\otimes L^{-1}                 \tag{DE-P13.1}
\]
has a filtration with four successive quotients \(\mathcal O_A\).
It is slope semistable of slope zero for every polarization.
Faithful flat pullback therefore implies semistability of \(F_*L\).

Suppose a saturated proper subsheaf \(Q\subset F_*L\), of
rank \(q\in\{1,2,3\}\), has equal slope. Put
\(S=F^*Q\otimes L^{-1}\subset U\) and \(T=U/S\).
Relative Frobenius is finite flat, so \(T\) is torsion free;
both \(S\) and \(T\) have slope zero.
Intersect \(S\) with the above \(\mathcal O_A\)-filtration.
Every nonzero successive quotient is a rank-one subsheaf of
\(\mathcal O_A\), hence is \(I_Z(-D)\) for an effective
divisor \(D\) and a zero-dimensional scheme \(Z\).
Their degrees are nonpositive and sum to zero. Thus each \(D=0\).
It follows that
\[
                   \det S\simeq\mathcal O_A,
                   \qquad \chi(S)\le0.                \tag{DE-P13.2}
\]
Dualizing gives \(T^*\subset U^*\). The same argument, applied
to the dual unipotent filtration and the slope-zero sheaf \(T^*\),
gives \(\chi(T^*)\le0\). On the smooth surface, \(T^{**}\)
is locally free and \(T^*=(T^{**})^*\). Riemann--Roch with
trivial canonical bundle gives
\[
 \chi(T)=\chi(T^{**})-\operatorname{length}(T^{**}/T)
        =\chi(T^*)-\operatorname{length}(T^{**}/T)\le0.
\]
Since \(\chi(U)=0\), additivity forces \(\chi(S)=0\).

Write \(l=c_1(L)\). From \(c_1(S)=0\) one has
\(F^*c_1(Q)=q l\). Riemann--Roch on an abelian surface,
\(\deg F=4\), and tensoring by \(L^{-1}\) now give
\[
 0=\chi(S)=4\chi(Q)-F^*c_1(Q)\cdot l+q\frac{l^2}{2}
          =4\chi(Q)-q\chi(L).                        \tag{DE-P13.3}
\]
Because \(\chi(L)\) is odd, this implies \(4\mid q\), a
contradiction. Thus no proper equal-slope subsheaf exists.
The surface stability assertion has been proved here; Sun's
curve stability theorem is not applied to a surface.
\end{proof}

\begin{lemma}[Three sections force a degree-six intersection]
\label{lem:DE-three-section-divisor-degree}
\status{Proved; full independent proof, source and code audits passed.}
Let \(H\) be an ample globally generated line bundle of type
\((1,3)\) on an abelian surface. If a line bundle \(D\) has
at least three independent sections, then
\[
                         c_1(D)\cdot c_1(H)\ge6.     \tag{DE-P13.4}
\]
\end{lemma}
\begin{proof}
Choose an effective divisor representing \(D\). Every effective
divisor on an abelian surface is nef: translate it to avoid
containing any given curve, and compute the nonnegative proper
intersection. If \(D^2>0\), the Hodge index theorem and the
absence of negative-square curves imply that \(D\) is ample.
Ample-line cohomology vanishing on an abelian variety and
Riemann--Roch give \(D^2=2h^0(D)\ge6\). Since \(H^2=6\),
Hodge index gives \((D\cdot H)^2\ge D^2H^2\ge36\), proving
the assertion in this case.

Suppose \(D^2=0\), and write \(D=\sum n_i C_i\) as an
effective divisor. All summands are nef, so \(C_i^2=0\) and
\(C_i\cdot C_j=0\) for every pair. Adjunction gives
\(p_a(C_i)=1\). An abelian variety contains no rational curve:
a map from \(\mathbf P^1\) factors through its zero Albanese.
Thus the normalization has genus one and its delta invariant is
zero. Each \(C_i\) is a smooth elliptic translate. After choosing
an origin, the map of that elliptic curve into the abelian surface
is a homomorphism, so its image is an elliptic subgroup translate.
Two different elliptic subgroup directions intersect positively:
their addition map onto the surface is an isogeny. The zero
intersections therefore force every \(C_i\) to be a fibre of
one elliptic quotient \(q:A\to B\).

Consequently \(D=q^*D_0\) for an effective divisor \(D_0\)
of degree \(n=\sum n_i\) on the elliptic curve \(B\).
Connected fibres give \(q_*\mathcal O_A=\mathcal O_B\), so
\(h^0(D)=h^0(D_0)=n\ge3\). On an elliptic fibre \(E\),
the restriction of \(H\) is globally generated and has positive
degree. Degree one is impossible, since its unique section has
a zero. Hence \(H\cdot E\ge2\), and
\(D\cdot H=n(E\cdot H)\ge6\).
This proves the degenerate-divisor case directly, in characteristic
two as well.
\end{proof}

\begin{lemma}[The primitive Cartier tensor is the nonzero alternating tensor]
\label{lem:DE-primitive-one-three-Cartier}
\status{Proved; full independent proof, source and code audits passed.}
Let \(A\) be nonordinary and let \(L\) be ample, globally
generated and of type \((1,3)\). Set \(W=H^0(A,L)\).
The relative Cartier pairing is an isomorphism
\[
 \beta:\bigwedge^2 W\longrightarrow H^0(A^{(2)},L^{(2)}).
                                                               \tag{DE-P13.5}
\]
In a level-three Schr\"odinger basis \(e_0,e_1,e_2\), and its
Frobenius-twisted target basis, there is \(c\in k^\times\) such that
\[
             \beta(e_i,e_j)=c e_k^{(2)}
             \quad\text{if }\{i,j,k\}=\{0,1,2\}.    \tag{DE-P13.6}
\]
\end{lemma}
\begin{proof}
Fix an invariant volume form on \(A\) and its Frobenius twist.
Finite-flat duality and projection formula give the perfect pairing
\[
          F_*L\otimes F_*L\longrightarrow L^{(2)}
                                                               \tag{DE-P13.7}
\]
defined by relative Cartier trace of products. Indeed
\(F^*L^{(2)}\simeq L^2\) and the relative dualizing line of
\(F\) is trivial after these volume trivializations. Nonordinarity
means that Cartier kills the invariant top form. Projection formula
therefore makes (DE-P13.7) alternating, including on local sections.
Its restriction to global sections is (DE-P13.5).

The odd theta group has the nondegenerate commutator pairing on
\(K(L)\simeq(\mathbf Z/3)^2\); see
\citep[Proposition~8.19 and Corollary~8.20]{EdixhovenVanDerGeerMoonenAV}.
Choose lifts of two generators with orders three and primitive
commutator \(\omega\). A clock eigenvector and its two shift
translates give a basis with shift \(e_i\mapsto e_{i+1}\) and
clock \(e_i\mapsto\omega^i e_i\). Distinct clock eigenvalues
and the transitive shift prove irreducibility. On \(\bigwedge^2W\)
the three clock eigenvalues are again distinct and the shift is
transitive, so this representation is irreducible of central weight
two. The target is the Frobenius twist of \(W\), also of weight
two. Relative Cartier commutes with translations and the line-bundle
isomorphisms defining the theta operators. Thus \(\beta\) is
equivariant and is zero or an isomorphism. The displayed clock
weights and shift force every nonzero equivariant map to have
exactly the form (DE-P13.6), with a single scalar.

It remains to rule out zero. Suppose \(\beta=0\), and let
\(Q\subset F_*L\) be the saturation of the global evaluation
image. Pairing (DE-P13.7) makes its generic fibre isotropic in
a four-dimensional symplectic space, so its rank \(q\) is one
or two. The saturation is theta invariant. If \(q=1\), its
determinant has at least the three independent sections of \(W\).
If \(q=2\), the wedge-evaluation map
\(\bigwedge^2W\to H^0(\det Q)\) is nonzero and equivariant.
The just-proved irreducibility makes it injective. Hence in both
cases \(h^0(\det Q)\ge3\).

On \(A^{(2)}\), write \(H=L^{(2)}\). One has numerically
\(c_1(F_*L)=2c_1(H)\), either from (DE-P13.1) or by
Riemann--Roch, so \(\mu_H(F_*L)=3\).
The stability lemma gives
\[
                     c_1(\det Q)\cdot H<3q\le6,
\]
contrary to (DE-P13.4). Therefore \(\beta\ne0\), proving all
assertions. No simplicity-to-nonvanishing inference is used.
\end{proof}

\begin{lemma}[A smooth member from principal descent]
\label{lem:DE-primitive-one-three-smooth-member}
\status{Proved; full independent proof, source and code audits passed.}
Every ample globally generated line bundle of type \((1,3)\)
on an abelian surface in characteristic two has a nonempty open
of smooth connected members in its complete linear system.
\end{lemma}
\begin{proof}
Choose an isotropic cyclic subgroup of order three in \(K(L)\).
The theta extension splits over this subgroup, giving an actual
linearization and descent along an etale isogeny
\(\eta:A\to B\) of degree three. The descended ample line
\(M\) has \(\chi(M)=\chi(L)/3=1\), so it is principal.
We use the same odd maximal-isotropic descent as in the earlier
polarization constructions, with an actual splitting of the theta
extension, not just invariance of the line-bundle class.

Here is the complete theta-divisor alternative used in this step.
Let \(\Theta\) be the effective divisor of the unique nonzero
section of \(M\); thus \(\Theta^2=2\). If its support is one
irreducible curve \(C\), write \(\Theta=nC\). Adjunction makes
\(C^2\) an even integer, and \(2=n^2C^2\) forces \(n=1\).
Then \(p_a(C)=2\). If \(C\) were singular, its normalization
would have genus at most one, and cannot have genus zero because
an abelian variety has no rational curves. A map from a genus-one
curve to an abelian variety is, after translation, a homomorphism.
Its image is a smooth elliptic translate. This contradicts both
singularity and \(p_a(C)=2\). Hence \(C\) is smooth of genus two.

If there are at least two components, write
\(\Theta=\sum n_i C_i\). Ampleness gives positive integers
\(\Theta\cdot C_i\), and
\[
                   2=\sum n_i(\Theta\cdot C_i).
\]
Therefore there are exactly two components, each of multiplicity
one, and \(\Theta\cdot C_1=\Theta\cdot C_2=1\).
All effective curves are nef, as above. Since each \(C_i^2\)
is a nonnegative even integer, these equalities force
\(C_1^2=C_2^2=0\) and \(C_1\cdot C_2=1\).
The preceding elliptic-curve argument makes both components
smooth elliptic translates. Intersection length one means that
they meet at exactly one point and transversely. This excludes
both additional components and hidden multiplicities. Translating
the meeting point to the origin, their addition map has degree
one and identifies this principal surface with their product,
although only the nodal assertion is needed here.
Thus in every case the etale pullback gives a member
\(C_0\in|L|\) that is smooth or reduced nodal. No nonproduct
principal-quotient existence assertion is required.

If \(C_0\) is smooth, openness already proves the assertion.
Otherwise choose \(s_1\in H^0(L)\) nonzero at every node of
\(C_0=(s_0=0)\). Global generation and the infinitude of \(k\)
allow this simultaneous avoidance of finitely many hyperplanes.
The pencil \(s_0+z s_1\) has a smooth generic fibre. To check
this without a separability version of Bertini, complete at a node.
Its equation is \(xy+z u(x,y,z)=0\), with \(u\) a unit.
On its relative singular scheme the partial derivatives imply
\(x,y\in(z)\). Substituting in the equation gives \(z\) times
a unit, so that singular scheme has no component dominating the
parameter trait. Away from the nodes the central fibre is already
smooth. Properness therefore excludes every generic singular point.
Finally every effective ample divisor here is connected: the
hypersurface sequence and \(H^1(A,L^{-1})=0\) give
\(H^0(C,\mathcal O_C)=k\). A smooth member is consequently
connected, and has genus \(1+L^2/2=4\).
\end{proof}

\begin{lemma}[The actual rank-two Euler map]
\label{lem:DE-primitive-one-three-Euler}
\status{Proved; full independent proof, source and code audits passed.}
Under the nonordinary hypotheses, let \(X=\mathbf P(W)\simeq
\mathbf P^2\). The quotient of the universal divisor's coherent
\(H^1\) by the constant ambient \(H^1(A,\mathcal O_A)\) is
\(\mathcal E=\Omega_X^1\), with regular linearized Frobenius
\(\Phi:F_X^*\mathcal E\to\mathcal E\). Its bijective locus is
\[
             X^\circ=X\setminus\{s_0^3+s_1^3+s_2^3=0\}.
                                                               \tag{DE-P13.8}
\]
On \(s_0=1\), writing \(t=s_1,u=s_2\), its matrix is
\[
 \Phi=c\begin{pmatrix}ut^2&1+u^3\\1+t^3&tu^2\end{pmatrix}.
                                                               \tag{DE-P13.9}
\]
The determinant has order one on the Fermat cubic. On its dense
open \(tu\ne0\), the first and stable ranks are one.
\end{lemma}
\begin{proof}
This assertion concerns every nonzero section, including a singular
zero divisor. The universal hypersurface sequence, ample-line
vanishing and Serre duality identify the quotient with
\[
 \ker\bigl(W^*\otimes\mathcal O_X(-1)\longrightarrow\mathcal O_X\bigr)
       =\Omega_X^1.                                   \tag{DE-P13.10}
\]
Frobenius preserves the constant ambient subbundle and induces a
regular map on this kernel. Duality with relative Cartier gives,
on a coefficient chart, the restriction of \(M_s^t\) from
\((s^{[2]})^\perp\) to \(s^\perp\), where
\(M_s(v)=\beta(s,v)\). This is the direct hypersurface calculation
of (OW12)--(OW13); it has no principal-cube hypothesis.

By (DE-P13.6),
\[
 M_s=c\begin{pmatrix}0&u&t\\u&0&1\\t&1&0\end{pmatrix}.
\]
Use the Euler frame \((ta+ub,a,b)\) in \(s^\perp\), and its
actual Frobenius pullback \((t^2a+u^2b,a,b)\) in
\((s^{[2]})^\perp\). Multiplying gives (DE-P13.9).
Its determinant is \(c^2(1+t^3+u^3)\). The analogous chart
calculations give the homogeneous equation (DE-P13.8), also
consistent with \(\det\Phi\in H^0(X,\mathcal O_X(3))\).
The Fermat cubic is smooth, since its three homogeneous partials
are \(s_0^2,s_1^2,s_2^2\); hence it is reduced and integral.

On this cubic the matrix is
\[
        c\binom{t^2}{u^2}(u\quad t).                 \tag{DE-P13.11}
\]
The scalar in its next semilinear iterate is, up to the nonzero
constant \(c^2\),
\(ut^4+tu^4=ut(t^3+u^3)=ut\). Thus the rank-one image is
Frobenius bijective when \(tu\ne0\), proving stable rank one.
Also every nonzero dual contraction of (DE-P13.6) is alternating
of rank two. All map and rank assertions have been checked on
the actual Euler bundle, not transferred from an unrelated odd
cover system.
\end{proof}

\begin{theorem}[Full relative image for a globally generated \((1,3)\) system]
\label{thm:DE-primitive-one-three-full-image}
\status{Proved; full independent proof, source and code audits passed.}
Let \(A\) be a fixed nonordinary abelian surface over \(k\), and
let \(L\) be ample, globally generated and of type \((1,3)\).
On the nonempty smooth-member, Frobenius-bijective relative open,
the rank-two relative Lang local system has geometric image
\[
                       \operatorname{GL}_2(\mathbf F_2).
                                                               \tag{DE-P13.12}
\]
Its nonzero cover is connected of degree three. A generic smooth
member has genus four and \(p\)-rank \(f(A)+2\).
\end{theorem}
\begin{proof}
First compute the image on the whole Euler bijective locus
\(X^\circ\) of (DE-P13.8). Its Lang fixed sheaf is finite etale
of rank two over \(\mathbf F_2\); the geometric image is a
subgroup of \(\operatorname{GL}_2(\mathbf F_2)\simeq S_3\).
In (DE-P13.9), replacing a Lang vector by \(c^{-1}\) times a
new vector removes \(c\) from its equation. Hence use \(c=1\).

Restrict first to the coefficient line \(t=0\), with
\(1+u^3\ne0\). The fixed-vector equations are
\[
                  a=(1+u^3)b^2,\qquad b=a^2.
\]
Its three nonzero vectors are parametrized by
\[
                         a^3=(1+u^3)^{-1}.           \tag{DE-P13.13}
\]
This Kummer cubic is irreducible over \(k(u)\): the rational
function on the right has simple poles at the three distinct
roots of \(1+u^3\), so it is not a cube. Since \(k\) contains
the third roots of unity, the line's image contains a three-cycle.

Restrict next to \(u=1\), with \(t\ne0\). The equations are
\[
             a=t^2a^2,\qquad b=(1+t^3)a^2+t b^2.
\]
There is one nonzero rational fixed vector \((0,t^{-1})\).
The other two have \(a=t^{-2}\); on writing \(w=tb\), their
equation is
\[
                          w^2+w=t^{-3}+1.            \tag{DE-P13.14}
\]
This is a nontrivial Artin--Schreier extension of \(k(t)\).
Indeed a pole of \(h^2+h\) has even order, whereas the right
side has a pole of order three at zero. Thus this line's image
contains a transposition of the three nonzero vectors. Both
line systems are actual restrictions of the same Euler system;
their images embed in the global image up to conjugacy.
A subgroup of \(S_3\) containing elements of orders three and
two is all of \(S_3\). This proves (DE-P13.12) on \(X^\circ\)
without invoking a higher-rank inertia or generation theorem.

The smooth-member locus is nonempty by the smoothing lemma.
It meets \(X^\circ\), since both are nonempty opens of the
integral plane. Restriction to this dense open leaves the finite
geometric image unchanged: a connected finite etale cover of a
normal integral base is integral and stays connected on every
nonempty open. The relative Lang system on the smooth family
is exactly the hypersurface-cohomology quotient (DE-P13.10),
so the computed image is its image. Neither line used above,
nor the Fermat wall, is required to parametrize smooth curves.

Finally, for a smooth member \(C\), the exact coherent sequence
\[
 0\longrightarrow H^1(A,\mathcal O_A)
  \longrightarrow H^1(C,\mathcal O_C)
  \longrightarrow\mathcal E_{[C]}\longrightarrow0
\]
is Frobenius compatible. Over the algebraically closed field the
stable fixed-vector functor is exact on finite-dimensional
Frobenius spaces: decompose the subspace into its bijective and
nilpotent parts, and solve the additive fixed-vector lifting
equation on each. Its dimensions therefore add. Since the last
space is bijective of dimension two on \(X^\circ\), the
generic \(p\)-rank is \(f(A)+2\). Genus four was proved above.
\end{proof}

\begin{remark}[Scope and source record]
\label{rem:DE-primitive-one-three-scope}
\status{Proved; full independent proof, source and code audits passed.}
This is a primitive-polarization source theorem under the explicit
global-generation hypothesis. It does not prove that every
polarization of type \((1,3)\) is globally generated, does not
give full image for arbitrary type \((1,\ell)\), and does not
settle the fixed principal degree-seventeen problem. No passage
from a generic moduli point to every fixed surface is used.

The only additional Frobenius-filtration source used is the full
higher-dimensional argument of Sun cited above. Its complete
Definition~3.1, Lemmas~3.2--3.5 and Theorem~3.6, including the
proofs on printed pages~5--9 of arXiv:math/0608043v2, were opened
and read. The odd-slope stability argument and the isotropic
effective-divisor classification are proved here. The smooth
member is obtained by principal descent and explicit nodal
smoothing, not an unqualified positive-characteristic Bertini
statement. The two explicit Lang slices prove full image directly.
The standalone exact program
\path{code/obstructions_crystals/verify_primitive_type_1_3.py}
checks 38 universal characteristic-two identities: the Euler-frame
multiplication, determinant and rank-one wall factorization, its
actual second semilinear iterate, constant scalar normalization,
both Lang slices, the rational fixed vector and the level-three
clock weights. The author run passed at 05:01:18.505300 UTC on
20 November 2026 in \(0.06905563100008294\) seconds. Exact
source, output and streams are retained. No finite sampling or
executed numerical failure is involved; the certificate does not
replace the geometric proofs.
\end{remark}
